\documentclass[a4paper,12pt]{paper}

\usepackage{mycommands}

% Worked-example layout: boxed problem statement followed by a solution.
\newcounter{example}
\renewenvironment{example}[1]{%  % paper.cls already defines an example environment
  \refstepcounter{example}%
  \begin{mdframed}[linewidth=1pt,innertopmargin=6pt,innerbottommargin=6pt]%
  \noindent\textbf{Example \theexample: #1.}\ }{%
  \end{mdframed}}
\newenvironment{solution}{\par\noindent\textbf{Solution.}\ }{\par\medskip}

\title{Worked examples for Lecture 14: Seismic sources}
\author{David Al-Attar \\ Michaelmas Term, \the\year}

\begin{document}

\maketitle

\noindent These examples accompany Lecture~14. They are intended as practice in the concrete
calculations that underlie the lecture, and are less demanding than the questions on the
problem sets. Numerical results were checked with the scripts in the \texttt{python}
directory.

\begin{example}{The stress glut of a fault in an isotropic body, and the moment tensor}
An idealised fault lies in the $(x_{1},x_{2})$-plane of an isotropic body with Lam\'{e}
parameters $\lambda$ and $\mu$, and the displacement discontinuity vector $\mathbf{w}$ is
tangential to the fault, so that $w_{3}=0$.
\begin{enumerate}
    \item[(a)] Starting from the stress glut $\overline{S}_{ij}=A_{ijk3}w_{k}\delta(x_{3})$
    obtained in the lecture, show that
    \begin{equation*}
    \overline{S}_{ij}=\mu\left(w_{i}\delta_{j3}+w_{j}\delta_{i3}\right)\delta(x_{3}),
    \end{equation*}
    and comment on the absence of $\lambda$.
    \item[(b)] The \textbf{moment tensor} of a source is defined as the volume integral of
    the stress glut,
    \begin{equation*}
    M_{ij}(t)=\int_{M}\overline{S}_{ij}(\mathbf{x},t)\dd^{3}\mathbf{x}.
    \end{equation*}
    Suppose that for $t>0$ the slip $\mathbf{w}=D\,\hat{\mathbf{e}}_{1}$ is uniform over a
    patch $\Sigma$ of the fault having area $\mathcal{A}$, that $\mu$ is constant on $\Sigma$,
    and that $\Sigma$ is small compared with the wavelengths of interest. Explain why the
    stress glut can then be replaced by $M_{ij}\delta(\mathbf{x})$, determine $M_{ij}$, and
    find its eigenvalues and eigenvectors. Express $\mathbf{M}$ in terms of its eigenvectors.
\end{enumerate}
\end{example}

\begin{solution}
(a) For an isotropic material the elastic tensor takes the form
\begin{equation}
A_{ijkl}=\lambda\delta_{ij}\delta_{kl}+\mu\left(\delta_{ik}\delta_{jl}+\delta_{il}\delta_{jk}\right),
\label{eq:1}
\end{equation}
and so, setting $l=3$ and contracting with $w_{k}$,
\begin{equation}
A_{ijk3}w_{k}=\lambda\delta_{ij}\delta_{k3}w_{k}+\mu\left(\delta_{ik}\delta_{j3}+\delta_{i3}\delta_{jk}\right)w_{k}
=\lambda w_{3}\,\delta_{ij}+\mu\left(w_{i}\delta_{j3}+w_{j}\delta_{i3}\right).
\label{eq:2}
\end{equation}
The first term is proportional to $w_{3}$, which vanishes because the displacement
discontinuity is tangential to the fault, and hence
\begin{equation}
\overline{S}_{ij}=\mu\left(w_{i}\delta_{j3}+w_{j}\delta_{i3}\right)\delta(x_{3}).
\label{eq:3}
\end{equation}
Written out, the only non-zero components are
\begin{equation}
\overline{S}_{13}=\overline{S}_{31}=\mu w_{1}\delta(x_{3}),\quad
\overline{S}_{23}=\overline{S}_{32}=\mu w_{2}\delta(x_{3}),
\label{eq:4}
\end{equation}
so that the stress glut is symmetric, as it must be within the linearised theory, and also
traceless, $\overline{S}_{ii}=2\mu w_{3}\delta(x_{3})=0$. The same calculation for a general
fault surface with unit normal $\hat{\mathbf{n}}$ gives
$\overline{S}_{ij}=\mu(w_{i}\hat{n}_{j}+w_{j}\hat{n}_{i})\delta_{\Sigma}$ whenever
$\mathbf{w}\cdot\hat{\mathbf{n}}=0$. The Lam\'{e} parameter $\lambda$ enters only through the
term $\lambda w_{3}\delta_{ij}$, which is to say only if the fault opens or the two sides
interpenetrate. For a shear dislocation the source is therefore completely insensitive to
$\lambda$; the radiated waves depend on $\lambda$ only through their propagation, not
through their generation.

(b) The stress glut enters the linearised equations of motion through
$-\partial\overline{S}_{ij}/\partial x_{j}$, and it acts on the displacement field only
through integrals against smooth functions. For any smooth $\phi$, using the defining
property of $\delta(x_{3})$ to reduce the volume integral to one over the fault,
\begin{equation}
\int_{M}\overline{S}_{ij}\,\phi\dd^{3}\mathbf{x}
=\int_{\Sigma}\mu\left(w_{i}\delta_{j3}+w_{j}\delta_{i3}\right)\phi\dd S
\approx\phi(\mathbf{0})\int_{\Sigma}\mu\left(w_{i}\delta_{j3}+w_{j}\delta_{i3}\right)\dd S
=\phi(\mathbf{0})M_{ij},
\label{eq:5}
\end{equation}
provided that $\phi$ varies little over $\Sigma$, which is the case for wavefields whose
wavelengths are large compared with the patch. Since
$\int_{M}M_{ij}\delta(\mathbf{x})\phi\dd^{3}\mathbf{x}=M_{ij}\phi(\mathbf{0})$, the
replacement $\overline{S}_{ij}\to M_{ij}\delta(\mathbf{x})$ preserves precisely this
zeroth moment of the glut, and this is the \textbf{point-source approximation}. Evaluating
the integral for uniform slip $\mathbf{w}=D\,\hat{\mathbf{e}}_{1}$ and constant $\mu$ gives,
for $t>0$,
\begin{equation}
M_{ij}=\mu\mathcal{A}D\left(\delta_{i1}\delta_{j3}+\delta_{i3}\delta_{j1}\right),
\quad\text{i.e.}\quad
\mathbf{M}=M_{0}\begin{pmatrix}0&0&1\\0&0&0\\1&0&0\end{pmatrix},\quad M_{0}=\mu\mathcal{A}D,
\label{eq:6}
\end{equation}
where $M_{0}$ is the \textbf{scalar seismic moment}, with units of $\mathrm{N\,m}$. A
moment tensor of this form is called a \textbf{double couple}, for reasons made clear in
the next example.

To find the eigenvalues we expand $\det(\mathbf{M}-\sigma\mathbf{I})$ along its second
row, which contains the single non-zero entry $-\sigma$, to obtain
\begin{equation}
\det(\mathbf{M}-\sigma\mathbf{I})=-\sigma\left(\sigma^{2}-M_{0}^{2}\right),
\label{eq:7}
\end{equation}
so the eigenvalues are $M_{0}$, $0$ and $-M_{0}$. The corresponding unit eigenvectors are
found by inspection, since $\mathbf{M}$ simply exchanges the first and third components of
a vector (scaled by $M_{0}$) and annihilates the second,
\begin{equation}
\hat{\mathbf{t}}=\frac{1}{\sqrt{2}}\left(\hat{\mathbf{e}}_{1}+\hat{\mathbf{e}}_{3}\right),\quad
\hat{\mathbf{b}}=\hat{\mathbf{e}}_{2},\quad
\hat{\mathbf{p}}=\frac{1}{\sqrt{2}}\left(\hat{\mathbf{e}}_{1}-\hat{\mathbf{e}}_{3}\right),
\label{eq:8}
\end{equation}
with $\mathbf{M}\hat{\mathbf{t}}=M_{0}\hat{\mathbf{t}}$, $\mathbf{M}\hat{\mathbf{b}}=\mathbf{0}$
and $\mathbf{M}\hat{\mathbf{p}}=-M_{0}\hat{\mathbf{p}}$. The vector $\hat{\mathbf{t}}$ lies at
$45^{\circ}$ to the fault plane in the direction of slip and is known as the
\textbf{tension} or T-axis; $\hat{\mathbf{p}}$, also at $45^{\circ}$ to the fault, is the
\textbf{pressure} or P-axis; and $\hat{\mathbf{b}}$, which lies in the fault plane
perpendicular to the slip, is the \textbf{null} or B-axis. Because $\mathbf{M}$ is
symmetric, it has the spectral decomposition
\begin{equation}
\mathbf{M}=M_{0}\left(\hat{\mathbf{t}}\hat{\mathbf{t}}^{T}-\hat{\mathbf{p}}\hat{\mathbf{p}}^{T}\right),
\label{eq:9}
\end{equation}
which can be checked directly since
$\hat{\mathbf{t}}\hat{\mathbf{t}}^{T}-\hat{\mathbf{p}}\hat{\mathbf{p}}^{T}
=\hat{\mathbf{e}}_{1}\hat{\mathbf{e}}_{3}^{T}+\hat{\mathbf{e}}_{3}\hat{\mathbf{e}}_{1}^{T}$.
Eq.~(\ref{eq:9}) says that a double couple is equivalent to a pair of equal and opposite
linear dipoles, a tensile one along $\hat{\mathbf{t}}$ and a compressive one along
$\hat{\mathbf{p}}$, at $45^{\circ}$ to the fault.

Two features of the result are worth noting. First, the trace of $\mathbf{M}$, which is
the sum of its eigenvalues, vanishes; a double couple involves no change in volume.
Second, $\mathbf{M}$ in eq.~(\ref{eq:6}) is unchanged if $\hat{\mathbf{e}}_{1}$ and
$\hat{\mathbf{e}}_{3}$ are interchanged, i.e.\ if the slip direction and the fault normal
swap roles. Slip $D\,\hat{\mathbf{e}}_{3}$ on a fault with normal $\hat{\mathbf{e}}_{1}$
produces an identical moment tensor. This second plane is known as the
\textbf{auxiliary plane}, and the moment tensor alone cannot distinguish it from the
true fault plane; both planes contain the null axis and bisect the P- and T-axes.
\end{solution}

\begin{example}{The equivalent body force for a point source}
A point source in the interior of a body has stress glut
$\overline{S}_{ij}(\mathbf{x},t)=M_{ij}(t)\delta(\mathbf{x})$, where the moment tensor
$M_{ij}$ is symmetric.
\begin{enumerate}
    \item[(a)] Show that the linearised equations of motion and boundary conditions are
    those of the same body subject to a body force
    $f_{i}=-M_{ij}\,\partial\delta(\mathbf{x})/\partial x_{j}$ and no applied surface
    traction.
    \item[(b)] Show that this force exerts no net force and no net torque on the body,
    \begin{equation*}
    \int_{M}f_{i}\dd^{3}\mathbf{x}=0,\quad
    \int_{M}\epsilon_{ijk}x_{j}f_{k}\dd^{3}\mathbf{x}=0,
    \end{equation*}
    and identify which property of $M_{ij}$ is responsible for each.
    \item[(c)] For the double couple of Example~1, interpret $\mathbf{f}$ as the limit of
    two force couples, and verify that their torques cancel.
\end{enumerate}
\end{example}

\begin{solution}
(a) The linearised equations of motion obtained in the lecture are
\begin{equation}
\rho\frac{\partial^{2}u_{i}}{\partial t^{2}}
-\frac{\partial}{\partial x_{j}}\left(A_{ijkl}\frac{\partial u_{k}}{\partial x_{l}}\right)
=-\frac{\partial\overline{S}_{ij}}{\partial x_{j}},
\label{eq:10}
\end{equation}
together with the boundary condition
$A_{ijkl}(\partial u_{k}/\partial x_{l})\hat{n}_{j}=\overline{S}_{ij}\hat{n}_{j}$ on
$\partial M$. Substituting $\overline{S}_{ij}=M_{ij}(t)\delta(\mathbf{x})$, and noting that
$M_{ij}$ does not depend on position, the right hand side of eq.~(\ref{eq:10}) becomes
\begin{equation}
f_{i}(\mathbf{x},t)=-M_{ij}(t)\frac{\partial\delta}{\partial x_{j}}(\mathbf{x}),
\label{eq:11}
\end{equation}
while $\delta(\mathbf{x})$ vanishes identically on $\partial M$ because the source lies in
the interior, so the boundary condition reduces to the traction-free condition
$A_{ijkl}(\partial u_{k}/\partial x_{l})\hat{n}_{j}=0$. The problem is therefore identical
to that for the same elastic body driven by the body force $\mathbf{f}$, which is why the
lecture refers to the right hand side of eq.~(\ref{eq:10}) as an effective body force.
The derivative of the delta function is understood in the sense of distributions: for any
smooth $\phi$,
\begin{equation}
\int_{M}\phi\,\frac{\partial\delta}{\partial x_{j}}\dd^{3}\mathbf{x}
=-\int_{M}\frac{\partial\phi}{\partial x_{j}}\,\delta\dd^{3}\mathbf{x}
=-\frac{\partial\phi}{\partial x_{j}}(\mathbf{0}),
\label{eq:12}
\end{equation}
where the surface term arising from integration by parts vanishes because $\delta$ is
supported only at the origin.

(b) Applying eq.~(\ref{eq:12}) with $\phi=1$ gives the net force
\begin{equation}
\int_{M}f_{i}\dd^{3}\mathbf{x}
=-M_{ij}\int_{M}\frac{\partial\delta}{\partial x_{j}}\dd^{3}\mathbf{x}
=M_{ij}\frac{\partial(1)}{\partial x_{j}}=0,
\label{eq:13}
\end{equation}
and this holds for any $M_{ij}$, symmetric or not; it is a consequence of the source
entering only as a divergence. For the torque about the origin we apply
eq.~(\ref{eq:12}) with $\phi=x_{j}$,
\begin{equation}
\int_{M}\epsilon_{ijk}x_{j}f_{k}\dd^{3}\mathbf{x}
=-\epsilon_{ijk}M_{kl}\int_{M}x_{j}\frac{\partial\delta}{\partial x_{l}}\dd^{3}\mathbf{x}
=\epsilon_{ijk}M_{kl}\,\frac{\partial x_{j}}{\partial x_{l}}\bigg|_{\mathbf{x}=\mathbf{0}}
=\epsilon_{ijk}M_{kl}\delta_{jl}
=\epsilon_{ijk}M_{kj}.
\label{eq:14}
\end{equation}
This is a contraction of the symmetric tensor $M_{kj}$ with the alternating symbol, which
is antisymmetric in $j$ and $k$, and so it vanishes. Since the net force is zero, the
torque about any other point is the same, and hence also zero. Conversely, if
$M_{ij}$ had an antisymmetric part, eq.~(\ref{eq:14}) shows that the source would exert a
net torque equal to minus twice the axial vector of that part. The symmetry of the stress
glut, which in the lecture was a consequence of the invariance of the strain energy under
superimposed rigid rotations, is thus exactly the condition that the source does not alter
the angular momentum of the body, while the divergence form guarantees that it does not
alter the linear momentum. An earthquake is an internal process, and cannot do either.

(c) For the double couple $M_{13}=M_{31}=M_{0}$ of Example~1, eq.~(\ref{eq:11}) gives
\begin{equation}
f_{1}=-M_{0}\frac{\partial\delta}{\partial x_{3}},\quad f_{2}=0,\quad
f_{3}=-M_{0}\frac{\partial\delta}{\partial x_{1}}.
\label{eq:15}
\end{equation}
A derivative of the delta function is the limit of a difference of two shifted delta
functions,
\begin{equation}
\frac{\partial\delta}{\partial x_{3}}(\mathbf{x})
=\lim_{h\to0}\frac{1}{h}\left[\delta\!\left(\mathbf{x}+\tfrac{h}{2}\hat{\mathbf{e}}_{3}\right)
-\delta\!\left(\mathbf{x}-\tfrac{h}{2}\hat{\mathbf{e}}_{3}\right)\right],
\label{eq:16}
\end{equation}
as follows by integrating both sides against a smooth $\phi$ and Taylor expanding. Hence
\begin{equation}
f_{1}=\lim_{h\to0}\frac{M_{0}}{h}\left[\delta\!\left(\mathbf{x}-\tfrac{h}{2}\hat{\mathbf{e}}_{3}\right)
-\delta\!\left(\mathbf{x}+\tfrac{h}{2}\hat{\mathbf{e}}_{3}\right)\right],
\label{eq:17}
\end{equation}
which is a point force of magnitude $F=M_{0}/h$ directed along $+\hat{\mathbf{e}}_{1}$ at
$\mathbf{x}=+\tfrac{h}{2}\hat{\mathbf{e}}_{3}$, together with an equal and opposite force
at $\mathbf{x}=-\tfrac{h}{2}\hat{\mathbf{e}}_{3}$. This is a \textbf{force couple} with
arm $h$ and moment $Fh=M_{0}$, which remains finite as $h\to0$; the forces act along the
slip direction and are separated along the fault normal. The side $x_{3}>0$ is pushed in
the $+\hat{\mathbf{e}}_{1}$ direction and the side $x_{3}<0$ in the $-\hat{\mathbf{e}}_{1}$
direction, consistent with the displacement jumping by $+D\,\hat{\mathbf{e}}_{1}$ on
crossing the fault from $x_{3}<0$ to $x_{3}>0$. Its torque is
\begin{equation}
\tfrac{h}{2}\hat{\mathbf{e}}_{3}\times F\hat{\mathbf{e}}_{1}
+\left(-\tfrac{h}{2}\hat{\mathbf{e}}_{3}\right)\times\left(-F\hat{\mathbf{e}}_{1}\right)
=Fh\,\hat{\mathbf{e}}_{3}\times\hat{\mathbf{e}}_{1}=M_{0}\hat{\mathbf{e}}_{2}.
\label{eq:18}
\end{equation}
In exactly the same way $f_{3}$ is a couple of forces $\pm F\hat{\mathbf{e}}_{3}$ at
$\pm\tfrac{h}{2}\hat{\mathbf{e}}_{1}$, with forces along the fault normal separated along
the slip direction, and torque
\begin{equation}
Fh\,\hat{\mathbf{e}}_{1}\times\hat{\mathbf{e}}_{3}=-M_{0}\hat{\mathbf{e}}_{2}.
\label{eq:19}
\end{equation}
The two torques cancel, as they must by part (b). The first couple is the one we might
have written down intuitively from the picture of slip on the fault; the second, which
corresponds to slip on the auxiliary plane, is forced upon us by the requirement that the
source exerts no net torque, and it is the pair together that constitutes the
\textbf{double couple}. A single couple would set the body spinning.
\end{solution}

\begin{example}{An explosive source generates only P-waves}
An explosion at the origin of a homogeneous isotropic whole space, with Lam\'{e}
parameters $\lambda$ and $\mu$ and density $\rho$, is modelled by the isotropic stress glut
\begin{equation*}
\overline{S}_{ij}(\mathbf{x},t)=M_{0}\,\delta_{ij}\,\delta(\mathbf{x})H(t).
\end{equation*}
Find the equivalent body force, write down the resulting equation of motion for
$\mathbf{u}$, and by taking its curl show that the displacement field is irrotational for
all time. Deduce that only P-waves are generated.
\end{example}

\begin{solution}
The equivalent body force is
\begin{equation}
f_{i}=-\frac{\partial}{\partial x_{j}}\left[M_{0}\delta_{ij}\delta(\mathbf{x})H(t)\right]
=-M_{0}H(t)\frac{\partial\delta}{\partial x_{i}}(\mathbf{x}),
\quad\text{i.e.}\quad
\mathbf{f}=-M_{0}H(t)\nabla\delta(\mathbf{x}),
\label{eq:20}
\end{equation}
which is the gradient of a scalar. For a homogeneous isotropic body we substitute
eq.~(\ref{eq:1}) into the elastic operator and use the constancy of $\lambda$ and $\mu$,
\begin{equation}
\frac{\partial}{\partial x_{j}}\left(A_{ijkl}\frac{\partial u_{k}}{\partial x_{l}}\right)
=\lambda\frac{\partial^{2}u_{k}}{\partial x_{i}\partial x_{k}}
+\mu\frac{\partial^{2}u_{i}}{\partial x_{j}\partial x_{j}}
+\mu\frac{\partial^{2}u_{j}}{\partial x_{j}\partial x_{i}}
=(\lambda+\mu)\frac{\partial}{\partial x_{i}}\left(\nabla\cdot\mathbf{u}\right)+\mu\nabla^{2}u_{i},
\label{eq:21}
\end{equation}
so that the equation of motion is the one given in Lecture~13 for a homogeneous isotropic
body, now with a body force,
\begin{equation}
\rho\frac{\partial^{2}\mathbf{u}}{\partial t^{2}}
-(\lambda+\mu)\nabla(\nabla\cdot\mathbf{u})-\mu\nabla^{2}\mathbf{u}
=-M_{0}H(t)\nabla\delta(\mathbf{x}),
\label{eq:22}
\end{equation}
to be solved with $\mathbf{u}=\partial\mathbf{u}/\partial t=\mathbf{0}$ at $t=0$. Now take
the curl of eq.~(\ref{eq:22}) and write $\bm{\omega}=\nabla\times\mathbf{u}$. The curl of
a gradient vanishes, which disposes of both the $(\lambda+\mu)$ term and the body force,
while the curl commutes with $\nabla^{2}$ and with $\partial^{2}/\partial t^{2}$ because the
coefficients are constant. We are left with
\begin{equation}
\rho\frac{\partial^{2}\bm{\omega}}{\partial t^{2}}-\mu\nabla^{2}\bm{\omega}=\mathbf{0},
\quad
\bm{\omega}(\mathbf{x},0)=\mathbf{0},\quad
\frac{\partial\bm{\omega}}{\partial t}(\mathbf{x},0)=\mathbf{0},
\label{eq:23}
\end{equation}
the initial conditions following by taking the curl of those for $\mathbf{u}$, which hold
at every point of the body. Each Cartesian component of $\bm{\omega}$ thus satisfies the
homogeneous scalar wave equation with speed $\beta=\sqrt{\mu/\rho}$ and vanishing initial
data, and the only such solution is zero. To see this, consider the energy
\begin{equation}
\mathcal{E}(t)=\frac{1}{2}\int\left(\rho\left\|\frac{\partial\bm{\omega}}{\partial t}\right\|^{2}
+\mu\frac{\partial\omega_{i}}{\partial x_{j}}\frac{\partial\omega_{i}}{\partial x_{j}}\right)\dd^{3}\mathbf{x},
\label{eq:24}
\end{equation}
whose time derivative, after an integration by parts, is
$\int\dot{\omega}_{i}\left(\rho\ddot{\omega}_{i}-\mu\nabla^{2}\omega_{i}\right)\dd^{3}\mathbf{x}=0$;
the surface term at infinity vanishes since, at any finite time, the disturbance has only
travelled a finite distance from the source. Thus $\mathcal{E}(t)=\mathcal{E}(0)=0$, and
as $\mathcal{E}$ is a sum of non-negative terms both $\partial\bm{\omega}/\partial t$ and
the gradient of $\bm{\omega}$ vanish; $\bm{\omega}$ is therefore constant in space and time,
and equal to its initial value $\mathbf{0}$. The displacement is \textbf{irrotational} at
all times.

In a whole space an irrotational field can be written as a gradient,
$\mathbf{u}=\nabla\phi$. Substituting into eq.~(\ref{eq:22}), and using
$\nabla^{2}\nabla\phi=\nabla\nabla^{2}\phi$ and $\nabla\cdot\nabla\phi=\nabla^{2}\phi$,
every term becomes a gradient,
\begin{equation}
\nabla\left[\rho\frac{\partial^{2}\phi}{\partial t^{2}}-(\lambda+2\mu)\nabla^{2}\phi
+M_{0}H(t)\delta(\mathbf{x})\right]=\mathbf{0}.
\label{eq:25}
\end{equation}
The bracket is therefore a function of $t$ alone, which can be absorbed into $\phi$
without changing $\mathbf{u}$, and so
\begin{equation}
\rho\frac{\partial^{2}\phi}{\partial t^{2}}-(\lambda+2\mu)\nabla^{2}\phi=-M_{0}H(t)\delta(\mathbf{x}).
\label{eq:26}
\end{equation}
This is the scalar wave equation with wave speed $\alpha=\sqrt{(\lambda+2\mu)/\rho}$, the
P-wave speed of Lecture~15: the explosion radiates a spherically symmetric P-wave and
nothing else. The isotropic moment tensor $M_{0}\delta_{ij}$ is invariant under all
rotations, so there is no fault plane and no preferred direction in the radiation. This
should be contrasted with the double couple of Example~1, whose moment tensor is
traceless. In practice explosions do produce S-waves, because the Earth is neither
homogeneous nor unbounded and P-waves are partly converted to S-waves at the free surface
and at internal discontinuities, but these are secondary. The size of the isotropic part
of the moment tensor, essentially zero for an earthquake, is one of the discriminants used
to distinguish nuclear tests from earthquakes.
\end{solution}

\begin{example}{Finite rupture and the moment rate function}
A rectangular fault occupies the region $0\le x_{1}\le L$, $0\le x_{2}\le W_{\mathrm{f}}$ of
the plane $x_{3}=0$ in a body with constant shear modulus $\mu$. Rupture begins at
$x_{1}=0$ at time $t=0$ and propagates in the $x_{1}$-direction at constant speed
$v_{\mathrm{r}}$, with the slip at each point taking place instantaneously as the rupture
front passes,
\begin{equation*}
w_{1}(x_{1},x_{2},t)=D\,H\!\left(t-\frac{x_{1}}{v_{\mathrm{r}}}\right),\quad w_{2}=w_{3}=0.
\end{equation*}
Find the time-dependent scalar moment $M_{0}(t)$ and the \textbf{moment rate}
$\dot{M}_{0}(t)$, and show that the latter is a boxcar of duration $L/v_{\mathrm{r}}$.
Estimate the rupture duration and the total moment for $L=100\,\mathrm{km}$,
$v_{\mathrm{r}}=3\,\mathrm{km\,s^{-1}}$, $W_{\mathrm{f}}=20\,\mathrm{km}$, $D=2\,\mathrm{m}$
and $\mu=3\times10^{10}\,\mathrm{Pa}$.
\end{example}

\begin{solution}
From eq.~(\ref{eq:4}) the stress glut has the single independent component
$\overline{S}_{13}=\overline{S}_{31}=\mu w_{1}\delta(x_{3})$, and so by the definition of
the moment tensor in Example~1 the only non-zero components are
\begin{equation}
M_{13}(t)=M_{31}(t)=\mu\int_{\Sigma}w_{1}\dd S\equiv M_{0}(t),
\label{eq:27}
\end{equation}
which defines the time-dependent scalar moment. The moment tensor has the double-couple
form of Example~1 at every instant, with the whole time dependence carried by
$M_{0}(t)$. Substituting the assumed slip, and performing the trivial integral over
$x_{2}$,
\begin{equation}
M_{0}(t)=\mu W_{\mathrm{f}}D\int_{0}^{L}H\!\left(t-\frac{x_{1}}{v_{\mathrm{r}}}\right)\dd x_{1}.
\label{eq:28}
\end{equation}
The step function equals one where $x_{1}<v_{\mathrm{r}}t$ and zero elsewhere, so the
integral is the length of the ruptured portion of the fault at time $t$, namely
\begin{equation}
M_{0}(t)=\begin{cases}
0, & t<0,\\
\mu W_{\mathrm{f}}D\,v_{\mathrm{r}}t, & 0\le t\le T_{\mathrm{r}},\\
\mu W_{\mathrm{f}}L D, & t>T_{\mathrm{r}},
\end{cases}
\quad T_{\mathrm{r}}=\frac{L}{v_{\mathrm{r}}}.
\label{eq:29}
\end{equation}
The moment grows linearly while the rupture front crosses the fault and then stays at its
final value $\mu\mathcal{A}D$, with $\mathcal{A}=LW_{\mathrm{f}}$, in agreement with
Example~1. Differentiating under the integral sign, and using $H'=\delta$,
\begin{equation}
\dot{M}_{0}(t)=\mu W_{\mathrm{f}}\int_{0}^{L}\frac{\partial w_{1}}{\partial t}\dd x_{1}
=\mu W_{\mathrm{f}}D\int_{0}^{L}\delta\!\left(t-\frac{x_{1}}{v_{\mathrm{r}}}\right)\dd x_{1}
=\mu W_{\mathrm{f}}D\,v_{\mathrm{r}}\int_{0}^{T_{\mathrm{r}}}\delta(t-\tau)\dd\tau,
\label{eq:30}
\end{equation}
where we substituted $\tau=x_{1}/v_{\mathrm{r}}$. The final integral is one if
$0<t<T_{\mathrm{r}}$ and zero otherwise, and so
\begin{equation}
\dot{M}_{0}(t)=\mu W_{\mathrm{f}}D\,v_{\mathrm{r}}\left[H(t)-H(t-T_{\mathrm{r}})\right],
\label{eq:31}
\end{equation}
a boxcar of height $\mu W_{\mathrm{f}}D\,v_{\mathrm{r}}$ and duration $T_{\mathrm{r}}$, in
agreement with the derivative of eq.~(\ref{eq:29}). Its area is
$\mu W_{\mathrm{f}}D\,v_{\mathrm{r}}T_{\mathrm{r}}=\mu\mathcal{A}D$, the total moment, as it
must be.

For the given values the rupture duration is
\begin{equation}
T_{\mathrm{r}}=\frac{100\,\mathrm{km}}{3\,\mathrm{km\,s^{-1}}}\approx33\,\mathrm{s},
\label{eq:32}
\end{equation}
and the total moment is
\begin{equation}
M_{0}=\mu\mathcal{A}D=\left(3\times10^{10}\,\mathrm{Pa}\right)
\left(10^{5}\,\mathrm{m}\right)\left(2\times10^{4}\,\mathrm{m}\right)\left(2\,\mathrm{m}\right)
=1.2\times10^{20}\,\mathrm{N\,m},
\label{eq:33}
\end{equation}
with a moment rate of $3.6\times10^{18}\,\mathrm{N\,m\,s^{-1}}$ during the rupture. For
orientation, the moment magnitude is defined by
$M_{\mathrm{w}}=\tfrac{2}{3}\left[\log_{10}(M_{0}/\mathrm{N\,m})-9.1\right]$, which here
gives $M_{\mathrm{w}}\approx7.3$, a large earthquake and a plausible size for a rupture of
these dimensions. The moment rate function, often called the source time function, is
what a distant seismometer effectively records: the far-field displacement radiated by a
point source is proportional to $\dot{M}_{0}$, so the pulse from this event would be
roughly $33\,\mathrm{s}$ long. Two refinements are easily incorporated. If the slip at each
point rises over a finite rise time $\tau_{\mathrm{r}}$ rather than instantaneously,
$\dot{M}_{0}$ becomes the convolution of the boxcar with a second boxcar of width
$\tau_{\mathrm{r}}$, i.e.\ a trapezoid of total duration $T_{\mathrm{r}}+\tau_{\mathrm{r}}$.
And for an observer in the direction of rupture propagation the pulse is compressed, as
successive parts of the fault radiate from progressively closer points, an effect known
as directivity. Note finally that the point-source approximation of Example~1 requires
wavelengths large compared with $L$, which for this event means periods well in excess of
$T_{\mathrm{r}}$.
\end{solution}

\newpage
\begin{example}{Well- and ill-posedness of polynomial interpolation}
Given distinct points $x_{1},\dots,x_{n}$ and values $y_{1},\dots,y_{n}$, we seek a
polynomial $p(x)=\sum_{k=0}^{m}c_{k}x^{k}$ with $p(x_{i})=y_{i}$. Write the problem as a
linear system $\mathbf{V}\mathbf{c}=\mathbf{y}$ for the coefficient vector $\mathbf{c}$,
and then:
\begin{enumerate}
    \item[(a)] for $n=3$ and $m=2$, show explicitly that a solution exists, is unique, and
    depends continuously on $\mathbf{y}$;
    \item[(b)] for $n=3$ and $m=1$, derive the condition that the data must satisfy for a
    solution to exist;
    \item[(c)] for $n=2$ and $m=2$, exhibit the one-parameter family of solutions;
    \item[(d)] for equally spaced points on $[0,1]$ with $m=n-1$, compute the condition
    number of $\mathbf{V}$ for $n=5$, $10$, $15$ and $20$, and comment on what continuous
    dependence is worth in practice.
\end{enumerate}
\end{example}

\begin{solution}
The interpolation conditions are linear in the unknown coefficients,
\begin{equation}
\sum_{k=0}^{m}x_{i}^{k}c_{k}=y_{i},\quad i=1,\dots,n,
\quad\text{i.e.}\quad
\mathbf{V}\mathbf{c}=\mathbf{y},\quad V_{ik}=x_{i}^{k},
\label{eq:34}
\end{equation}
where $\mathbf{V}$ is the $n\times(m+1)$ \textbf{Vandermonde matrix}. In the language of
the lecture the inputs are $\mathbf{y}$ and the outputs are $\mathbf{c}$. A solution exists
if and only if $\mathbf{y}$ lies in the column space of $\mathbf{V}$, it is unique if and
only if the null space of $\mathbf{V}$ is trivial, and, whenever the solution is unique,
it depends continuously on $\mathbf{y}$ because it is then given by a linear map between
finite-dimensional spaces, and every such map is continuous.

(a) With $n=3$ and $m=2$ the matrix is square, and subtracting the first row from the
other two,
\begin{equation}
\det\mathbf{V}=\det\begin{pmatrix}1&x_{1}&x_{1}^{2}\\1&x_{2}&x_{2}^{2}\\1&x_{3}&x_{3}^{2}\end{pmatrix}
=\det\begin{pmatrix}x_{2}-x_{1}&x_{2}^{2}-x_{1}^{2}\\x_{3}-x_{1}&x_{3}^{2}-x_{1}^{2}\end{pmatrix}
=(x_{2}-x_{1})(x_{3}-x_{1})(x_{3}-x_{2}),
\label{eq:35}
\end{equation}
where the last step factors $x_{j}^{2}-x_{1}^{2}=(x_{j}-x_{1})(x_{j}+x_{1})$ from each row.
This is non-zero for distinct points, so $\mathbf{V}$ is invertible and the problem has
the unique solution $\mathbf{c}=\mathbf{V}^{-1}\mathbf{y}$ for every $\mathbf{y}$. The
solution can also be written down directly in Lagrange form,
\begin{equation}
p(x)=\sum_{i=1}^{3}y_{i}\,\ell_{i}(x),\quad
\ell_{i}(x)=\prod_{j\neq i}\frac{x-x_{j}}{x_{i}-x_{j}},
\label{eq:36}
\end{equation}
since $\ell_{i}(x_{j})=\delta_{ij}$; this establishes existence, and uniqueness follows
because the difference of two solutions is a quadratic with three distinct roots, hence
identically zero. Continuous dependence is quantified by
\begin{equation}
\|\delta\mathbf{c}\|\le\|\mathbf{V}^{-1}\|\,\|\delta\mathbf{y}\|,
\quad
\frac{\|\delta\mathbf{c}\|}{\|\mathbf{c}\|}\le\mathrm{cond}(\mathbf{V})\frac{\|\delta\mathbf{y}\|}{\|\mathbf{y}\|},
\quad
\mathrm{cond}(\mathbf{V})=\|\mathbf{V}\|\,\|\mathbf{V}^{-1}\|,
\label{eq:37}
\end{equation}
where $\|\cdot\|$ is the matrix norm induced by the Euclidean vector norm and
$\mathrm{cond}(\mathbf{V})$ is the \textbf{condition number}, which for this norm is the
ratio of the largest to the smallest singular value of $\mathbf{V}$. The second inequality
follows from the first together with $\|\mathbf{y}\|\le\|\mathbf{V}\|\,\|\mathbf{c}\|$.
The problem is well-posed.

(b) With $n=3$ and $m=1$ we have three equations $c_{0}+c_{1}x_{i}=y_{i}$ for two
unknowns. The two columns of $\mathbf{V}$, namely $(1,1,1)^{T}$ and $(x_{1},x_{2},x_{3})^{T}$,
are linearly independent, so the column space is a plane in $\mathbb{R}^{3}$ and
$\mathbf{y}$ must lie in it. The plane is the orthogonal complement of any non-zero vector
$\mathbf{a}$ with $\mathbf{a}^{T}\mathbf{V}=\mathbf{0}$, and
\begin{equation}
\mathbf{a}=\left(x_{2}-x_{3},\;x_{3}-x_{1},\;x_{1}-x_{2}\right)^{T}
\label{eq:38}
\end{equation}
does the job, since its components sum to zero and
$(x_{2}-x_{3})x_{1}+(x_{3}-x_{1})x_{2}+(x_{1}-x_{2})x_{3}=0$. Existence therefore requires
\begin{equation}
(x_{2}-x_{3})y_{1}+(x_{3}-x_{1})y_{2}+(x_{1}-x_{2})y_{3}=0,
\label{eq:39}
\end{equation}
and this condition is also sufficient. Rearranging, it reads
$(y_{3}-y_{1})(x_{2}-x_{1})=(y_{2}-y_{1})(x_{3}-x_{1})$, which states that the slope from
the first point to the third equals the slope from the first to the second: the three
points $(x_{i},y_{i})$ must be collinear. Equivalently, the determinant of the matrix with
rows $(1,x_{i},y_{i})$ vanishes. For equally spaced points $x_{i}=x_{1}+(i-1)h$ the
condition is $y_{1}-2y_{2}+y_{3}=0$, the vanishing of the second difference. When it holds,
the solution is $c_{1}=(y_{2}-y_{1})/(x_{2}-x_{1})$ and $c_{0}=y_{1}-c_{1}x_{1}$. Data
chosen at random satisfy eq.~(\ref{eq:39}) with probability zero, so for noisy
measurements the problem as posed has no solution, and one is led, as noted in the lecture,
to relax what is meant by a solution, for instance to the $\mathbf{c}$ that minimises
$\|\mathbf{V}\mathbf{c}-\mathbf{y}\|$.

(c) With $n=2$ and $m=2$ there are two equations for three unknowns. A vector $\mathbf{c}$
lies in the null space if and only if the corresponding quadratic vanishes at $x_{1}$ and
$x_{2}$, that is, if and only if $p(x)=\alpha(x-x_{1})(x-x_{2})$ for some $\alpha$, so the
null space is spanned by
\begin{equation}
\mathbf{c}_{0}=\left(x_{1}x_{2},\;-(x_{1}+x_{2}),\;1\right)^{T},
\label{eq:40}
\end{equation}
as may be checked directly: $x_{1}x_{2}-(x_{1}+x_{2})x_{i}+x_{i}^{2}=(x_{i}-x_{1})(x_{i}-x_{2})=0$.
A particular solution is the straight line through the two points, and the general
solution is
\begin{equation}
p(x)=y_{1}+\frac{y_{2}-y_{1}}{x_{2}-x_{1}}(x-x_{1})+\alpha(x-x_{1})(x-x_{2}),
\label{eq:41}
\end{equation}
or in terms of coefficients, with $s=(y_{2}-y_{1})/(x_{2}-x_{1})$,
\begin{equation}
c_{2}=\alpha,\quad c_{1}=s-\alpha(x_{1}+x_{2}),\quad c_{0}=y_{1}-sx_{1}+\alpha x_{1}x_{2},
\label{eq:42}
\end{equation}
for arbitrary $\alpha$. Every quadratic through the two data points fits them equally
well, and the data provide no information whatever about $\alpha$. This is the
underdetermined case of the lecture, and it is the situation in which almost all
geophysical inverse problems find themselves. To select one member of the family we must
add information not contained in the data, for example by choosing the $\alpha$ that
minimises $\|\mathbf{c}\|$, and later lectures will have much to say about how such
choices are made and justified.

(d) For $m=n-1$ the matrix is square and, by the argument of part (a) generalised to $n$
points, invertible, so the problem is well-posed for every $n$. Taking equally spaced
points $x_{i}=(i-1)/(n-1)$ on $[0,1]$ and computing the condition number in the Euclidean
norm numerically,\footnote{For $n=20$ the condition number is comparable with the
reciprocal of double-precision machine epsilon, and a floating-point computation of it
is itself uncertain at the level of a few per cent; a routine such as
\texttt{numpy.linalg.cond} returns $1.2\times10^{16}$. The value quoted was obtained by
inverting $\mathbf{V}$ in exact rational arithmetic before taking norms.} we find
\begin{equation}
\begin{array}{c|cccc}
n & 5 & 10 & 15 & 20\\ \hline
\mathrm{cond}(\mathbf{V}) & 6.9\times10^{2} & 1.5\times10^{7} & 4.0\times10^{11} & 1.1\times10^{16}
\end{array}\,.
\label{eq:43}
\end{equation}
The growth is exponential in $n$. By eq.~(\ref{eq:37}) the condition number bounds the
factor by which relative errors in the data can be amplified in the coefficients, and the
bound is attained for suitably chosen perturbations. For $n=10$, data accurate to one part
in $10^{7}$ can produce coefficients with errors of order one hundred per cent; for $n=20$,
data accurate to sixteen significant figures, i.e.\ to the full precision of double
arithmetic, can produce coefficients with no correct digits at all. The problem is
well-posed in the strict sense of the lecture, since the solution does depend continuously
on the data, but the constant in that continuity is so large that the solution is
practically useless. Once measurement errors and finite precision are admitted, the
distinction between an ill-posed problem and a badly conditioned one all but disappears,
and both call for the same kind of remedy. It is worth adding that the ill-conditioning
here lies in the choice of unknowns: the monomials $x^{k}$ on $[0,1]$ become nearly
parallel as $k$ grows, so the columns of $\mathbf{V}$ are nearly linearly dependent. The
values of the interpolating polynomial are far better determined than its monomial
coefficients, and a wiser parameterisation, for instance in terms of Chebyshev
polynomials, greatly reduces the condition number. Choosing the model parameters well is
often the first step in taming an inverse problem.
\end{solution}

\end{document}
